\section{Experiments}
\label{sec:experiments}

We organise the evaluation around four claims. First, we show through analytical bridges that the cylindrical target has no origin singularity. Second, we verify the practical interaction of this geometry with U-Net architectures and Optimal Transport. Third, we ablate the framework across spatial dimensions to test scaling robustness and the effect of the joint OT coupling. Fourth, we show that factorising the coupling destroys the joint target distribution.

\subsection{Analytical Bridge Ablation: Isolating the Geometry}
\label{sec:exp_ablation}

To isolate the effect of the metric without the confounding variables of neural network capacity or imperfect empirical Optimal Transport coupling, we first evaluate the exact analytical probability paths. By computing the theoretical trajectories $z_t$ between samples drawn from a uniform phase prior and a complex target, we can directly measure the angular velocities predicted in Section~\ref{sec:method}. We run the same measurement against three targets: the synthetic copula used throughout this section, the coil-combined knee MRI of Section~\ref{sec:exp_unet_baseline}, and complex STFT segments of read speech. No network is involved in any of them. Because the numerator of the Cartesian angular velocity $\dot{\theta}(t) = \text{Im}(\bar{z}_0 z_1) / |z_t|^2$ is invariant along the straight-line path, the peak angular velocity is governed exclusively by the minimum distance to the origin.

Table~\ref{tab:bridge_ablation} reports the empirical distribution of peak angular velocities over $10^6$ sampled endpoint pairs (seed $0$); the peak on each bridge is evaluated in closed form (Proposition~\ref{thm:angular_divergence}), so the trajectories are exact and only the endpoint sampling is Monte Carlo. The Cartesian bridge exhibits a heavy-tailed power law (index $\approx 1.0$), consistent with Corollary~\ref{cor:pareto_tail}. With independent coupling, $46\%$ of Cartesian paths exceed an angular velocity of $\pi$ rad per unit time, and the 99.9th percentile reaches $1{,}566$; because the tail index is close to $1$, the sample maximum grows with the number of paths and is not a property of the bridge. Minibatch OT on scalar pairs partially mitigates this by avoiding origin-crossing pairs; the transport-cost reduction that makes this possible shrinks with field dimension (Section~\ref{sec:exp_scaling}), and we do not measure the bridges at field scale. In stark contrast, the Cylindrical bridge is structurally bounded by $\pi$ by construction ($0.0\%$ paths $> \pi$), independently of the coupling. The two real targets agree with the synthetic one on the size of the effect and are worse than it: $56.4\%$ of Cartesian paths exceed $\pi$ on knee MRI and $94.8\%$ on speech. The second column weights each path by its target's $|z_1|^2$ instead of counting coefficients equally, and the two disagree sharply on speech ($94.8\%$ against $49.4\%$). That disagreement is not an uncertainty in the measurement but the same modelling choice this paper makes explicit: weighting an angular quantity by $|z_1|^2$ reinstates exactly the $A^2$ that the plane's inherited metric $\mathrm{d}A^2 + A^2\mathrm{d}\theta^2$ carries on its angular term and that the cylinder drops. We therefore report both. This bounded regression target is consistent with the few-step results below.

\begin{table}[t]
    \centering
    \caption{\textbf{Analytical path geometry, on three sources of complex data.} Peak angular velocity $|\dot{\theta}|$ along $10^6$ exact bridges from the prior to each target (no network). \textbf{Paths} counts every coefficient equally; \textbf{Energy} weights each path by its target's $|z_1|^2$. Weighting by $|z_1|^2$ reinstates the $A^2$ factor from the Euclidean metric $\mathrm{d}A^2 + A^2\mathrm{d}\theta^2$ that the cylinder drops; we report both. \textbf{Targets.} \emph{Synthetic} copula field (Section~\ref{sec:exp_unet_baseline}); \emph{Knee MRI} fastMRI CORPD ($320\times320$); \emph{Speech STFT} LibriSpeech ($64\times64$ bins/frames). Normalised by each field's peak modulus.}
    \label{tab:bridge_ablation}
    \vspace{1mm}
    \begin{tabular}{llcc}
        \toprule
        \textbf{Target} & \textbf{Geometry / coupling} & \textbf{Paths $> \pi$ $\downarrow$} & \textbf{Energy $> \pi$ $\downarrow$} \\
        \midrule
        Synthetic
          & Cartesian, independent  & 46.0\% & 42.6\% \\
          & Cartesian, minibatch OT & 1.0\%  & 0.1\%  \\
          & \textbf{Cylindrical (ours)} & \textbf{0.0\%} & \textbf{0.0\%} \\
        \midrule
        Knee MRI
          & Cartesian, independent  & 56.4\% & 48.3\% \\
          & Cartesian, minibatch OT & 13.1\% & 0.1\%  \\
          & \textbf{Cylindrical (ours)} & \textbf{0.0\%} & \textbf{0.0\%} \\
        \midrule
        Speech STFT
          & Cartesian, independent  & 94.8\% & 49.4\% \\
          & Cartesian, minibatch OT & 81.9\% & 5.1\%  \\
          & \textbf{Cylindrical (ours)} & \textbf{0.0\%} & \textbf{0.0\%} \\
        \bottomrule
    \end{tabular}
    \vspace{1mm}
\end{table}


\subsection{Spatial Field Synthesis: U-Net Evaluation}
\label{sec:exp_unet_baseline}

While the analytical bridges characterise the two regression targets exactly, practical generative modeling relies on a neural network's ability to regress this vector field.

\textbf{Data and protocol.} All spatial experiments use a synthetic complex field with a known ground truth: amplitude and phase are coupled through a Gaussian copula (correlation $\rho = 0.5$) with fixed marginals, and the latents are spatially smoothed (correlation length 4 pixels, lag-1 amplitude correlation $0.96$). Each field is normalised by its own peak modulus, and metrics are computed in this training domain. We train a convolutional U-Net for 40 epochs (batch 64, 8192 training fields, the squared velocity error of Section~\ref{sec:method} in both geometries) with both geometries sharing the same prior, amplitude $\mathcal{U}[0,1]$ and phase $\mathcal{U}[0, 2\pi)$ independently per pixel (supported on the unit disc, though not uniform on it), and report the mean over 5 seeds. We call a difference separated when every seed of one arm lies below every seed of the other; for 5 against 5 seeds this is an exact two-sided Mann--Whitney $p = 0.008$, the smallest attainable. The headline claim is an intersection, a lower error at every evaluated step count and resolution, so every component has to pass on its own and the conjunction needs no multiplicity correction. With 5 seeds per arm the tests have little power, so where we report no significant difference that is a failure to detect one, not evidence of equivalence; an equivalence claim would need a prespecified margin and a test against it. Generation uses a Heun solver; $k$ steps correspond to $2k-1$ function evaluations. The generative error is the sliced $W_2$ distance between the per-pixel complex values of generated and reference fields, pooled over pixels. It compares pixel-level statistics and is insensitive to spatial arrangement (Section~\ref{sec:exp_factorized}), so the tables below do not measure spatial structure. We compare Cartesian and Cylindrical geometries under both Independent and joint OT couplings (Section~\ref{sec:ot_coupling}; for the Cartesian arm the same assignment uses the Euclidean cost).

Table~\ref{tab:unet_baseline} presents the results at $64\times64$. CyFM with the joint coupling has the lowest error at every $k \le 8$, with all seeds separated from the best Cartesian arm. At a single integration step ($k=1$) it reaches $0.117$, a $3\times$ lower error than the best Cartesian arm ($0.376$). At convergence ($k=100$) we detect no significant difference ($0.040$ against $0.046$, $p = 0.15$). With the joint coupling the cylindrical error decreases monotonically with $k$, here and at $32\times32$ (Table~\ref{tab:unet_scaling}); without it, the error rises from $k=1$ to $k=2$ at every resolution. Straightness is similar across arms ($0.13$--$0.14$); we did not test for equality. The two error components behave differently: at $k \le 8$ CyFM has the lower amplitude $W_2$ at every resolution with separated seeds, while the circular phase $W_2$ is mixed, the Cartesian geometry being lower at $k = 8$ on $32\times32$ and $64\times64$ fields. We report this as a decomposition of where the arms differ, not as evidence that the amplitude component causes the aggregate difference.

\begin{table}[h]
    \centering
    \caption{\textbf{Few-step synthesis, two domains.} Sliced $W_2$ between generated and
    reference complex fields, mean over 5 seeds, evaluated on held-out data. $k$ is the
    number of solver steps; a Heun step costs two function evaluations. $k^\star$ is the
    first step count at which an arm reaches $E^\star$, the $k=100$ error of the
    \emph{worse} geometry, which is what stops the comparison from being an artefact of
    the two converging to different asymptotes. Errors are comparable within a block and
    not across blocks. Each block holds its own protocol fixed across arms -- same architecture, batch, objective, seeds and epoch count --
    but an epoch is not the same amount of optimisation in the two domains: the speech store
    holds $110{,}099$ segments against $8{,}192$ synthetic fields, so forty epochs are
    thirteen times as many gradient steps there. Passes were held fixed, not steps; we say so
    because more data on one side is not a neutral choice in a comparison between domains.}
    \label{tab:unet_baseline}
    \vspace{1mm}
    \resizebox{\linewidth}{!}{
    \begin{tabular}{llccccc c}
        \toprule
        \textbf{Domain} & \textbf{Arm} & $k=1$ & $k=2$ & $k=4$ & $k=8$ & $k=100$ & $\mathbf{k^\star}\ \downarrow$ \\
        \midrule
        Synthetic $64\times64$
          & \textbf{Cylindrical + OT (ours)} & \textbf{0.1174} & \textbf{0.1141} & \textbf{0.0777} & \textbf{0.0537} & \textbf{0.0399} & \textbf{11} \\
          & Cylindrical, independent         & 0.1214 & 0.1478 & 0.1087 & 0.0665 & 0.0434 & 14 \\
          & Cartesian + OT                   & 0.3763 & 0.3180 & 0.2600 & 0.1409 & 0.0468 & 24 \\
          & Cartesian, independent           & 0.4159 & 0.3494 & 0.2849 & 0.1538 & 0.0458 & 32 \\
        \midrule
        Speech STFT $64\times64$
          & \textbf{Cylindrical + OT (ours)} & \textbf{0.0454} & \textbf{0.0403} & 0.0333 & 0.0308 & 0.0264 & \textbf{4} \\
          & Cylindrical, independent         & 0.0458 & 0.0409 & 0.0336 & \textbf{0.0303} & \textbf{0.0261} & \textbf{4} \\
          & Cartesian + OT                   & 0.0780 & 0.0695 & 0.0563 & 0.0442 & 0.0353 & 100 \\
          & Cartesian, independent           & 0.0822 & 0.0758 & 0.0610 & 0.0457 & 0.0353 & 100 \\
        \bottomrule
    \end{tabular}
    }
    \vspace{1mm}
\end{table}


\subsection{Dimensionality Scaling and the Joint OT Coupling}
\label{sec:exp_scaling}

The transport-cost reduction that minibatch OT achieves over independent pairing collapses with dimension. Under the cylindrical metric on the pointwise copula target (batch 64), it is $86\%$ for scalar pairs, $21\%$ for $8\times8$ fields and $3\%$ for $64\times64$ fields, while the assignment still reorders nearly every pair. To test whether the trained benefit collapses with it, we ablate the U-Net models across dimensions ($16\times16$, $32\times32$, and $64\times64$), with the cylindrical model trained both with and without the joint coupling (Table~\ref{tab:unet_scaling}).

It does not. At every resolution and every step count, joint cylindrical OT lowers the mean error of the cylindrical model; at $k \le 4$ the reduction is $3$--$60\%$, and the five seeds are separated in 6 of the 9 cases (the exceptions are $k=1$ at $16\times16$ and $k=2, 4$ at $64\times64$). With the coupling, CyFM has a lower error than the best Cartesian arm at every $k \le 8$ at all three resolutions, with all seeds separated, and at $k=100$ no difference is significant. Giving each geometry its best objective from $\{\mathcal{L}_1, \mathcal{L}_2\} \times \{\text{independent}, \text{OT}\}$ leaves this unchanged except at $16\times16$, $k=8$ ($p = 0.15$); at $64\times64$, $k=100$, the best arms reach $0.035$ and $0.032$ ($p = 0.69$; Appendix Table~\ref{tab:unet_loss_protocols}). OT is also most effective for the Cartesian geometry at low resolution: at $16\times16$, Cartesian $k=1$ improves from $0.442$ without OT to $0.227$ with it.

\begin{table}[t]
    \centering
    \caption{\textbf{Dimensionality Scaling.} Generative error (Sliced $W_2$, $\downarrow$) across spatial resolutions, mean over 5 seeds, with the cylindrical model trained with and without the joint OT coupling. All arms use the unweighted squared velocity error. CyFM's few-step advantage holds at all three resolutions, and joint OT improves the cylindrical model at every step count.}
    \label{tab:unet_scaling}
    \vspace{1mm}
    \resizebox{\linewidth}{!}{
    \begin{tabular}{lllccccc}
        \toprule
        \textbf{Resolution} & \textbf{Geometry} & \textbf{Coupling} & \textbf{1 Step} & \textbf{2 Steps} & \textbf{4 Steps} & \textbf{8 Steps} & \textbf{100 Steps} \\
        \midrule
        $16\times16$ ($D=256$)
        & Cartesian & OT & 0.227 & 0.202 & 0.153 & 0.097 & 0.068 \\
        & CyFM & Independent & 0.116 & 0.155 & 0.133 & 0.097 & 0.062 \\
        & \textbf{CyFM (Ours)} & \textbf{Joint OT} & \textbf{0.098} & \textbf{0.086} & \textbf{0.053} & \textbf{0.064} & 0.054 \\
        \midrule
        $32\times32$ ($D=1024$)
        & Cartesian & OT & 0.340 & 0.275 & 0.209 & 0.104 & 0.037 \\
        & CyFM & Independent & 0.123 & 0.145 & 0.113 & 0.063 & 0.047 \\
        & \textbf{CyFM (Ours)} & \textbf{Joint OT} & \textbf{0.113} & \textbf{0.094} & \textbf{0.072} & \textbf{0.049} & 0.041 \\
        \midrule
        $64\times64$ ($D=4096$)
        & Cartesian & OT & 0.376 & 0.318 & 0.260 & 0.141 & 0.047 \\
        & CyFM & Independent & 0.121 & 0.148 & 0.109 & 0.066 & 0.043 \\
        & \textbf{CyFM (Ours)} & \textbf{Joint OT} & \textbf{0.117} & \textbf{0.114} & \textbf{0.078} & \textbf{0.054} & 0.040 \\
        \bottomrule
    \end{tabular}
    }
    \vspace{1mm}
\end{table}


\subsection{The Factorized Coupling Trap}
\label{sec:exp_factorized}

Because the cylindrical metric is a sum over amplitude and phase, and over pixels, it is tempting to solve the transport problem separately per coordinate or per spatial patch. The factorised minimum is always a lower bound on the joint one, and it is attained exactly when some joint plan is simultaneously optimal in every coordinate --- which holds when both distributions are product measures, and degenerately when source and target coincide, but not in general. We measure the consequence without any network, on the pointwise copula target (spiral structure) with $n = 1024$ pairs and 8 seeds (Table~\ref{tab:factorized}). Coupling amplitude and phase independently leaves both marginals bit-identical to the target (Kolmogorov--Smirnov distance $0$), yet the circular-linear correlation of the coupled data falls to $\approx 0.04$ at every $\rho$, and its transport cost stays flat at $0.20$ while the true optimum rises to $0.38$ at $\rho = 1$: the factorised plan reports a cost below the optimum because it no longer transports the target. The spatial analogue behaves identically. On $64\times64$ fields (correlation length 6), coupling $16\times16$ patches independently leaves the pooled sliced $W_2$ to the data at $0.000$ and the correlation inside patches at $0.979$, but the lag-1 correlation across patch seams falls from $0.979$ to $0.044$. The damage is in the coupled targets themselves: a factorised plan reassembles endpoints from different data samples, so the law that the flow-matching objective transports to at its population optimum is that reassembled law and not the data. Capacity does not change which target is being fitted, and pooled distributional metrics do not detect the difference.

\begin{table}[h]
    \centering
    \caption{\textbf{The Factorized Coupling Trap} (no network). Coupling amplitude and phase independently preserves both marginals exactly but erases their dependence, and reports a transport cost below the true optimum. Spiral copula target, $n = 1024$, mean over 8 seeds.}
    \label{tab:factorized}
    \vspace{1mm}
    \begin{tabular}{cccccc}
        \toprule
        $\rho$ & \textbf{Corr.\ target} & \textbf{Corr.\ coupled} & \textbf{KS (marginals)} & \textbf{Cost, factorised} & \textbf{Cost, joint OT} \\
        \midrule
        0.0 & 0.024 & 0.042 & 0 & 0.200 & 0.206 \\
        0.5 & 0.310 & 0.042 & 0 & 0.200 & 0.225 \\
        1.0 & 0.867 & 0.040 & 0 & 0.200 & 0.385 \\
        \bottomrule
    \end{tabular}
    \vspace{1mm}
\end{table}

\subsection{Scope and Future Work}
\label{sec:exp_downstream}

All results in this paper are on synthetic complex fields. The next steps are: evaluating on physical signals, namely unconditional generation of complex MRI images (fastMRI, \citealp{zbontar2018fastmri}) and complex audio spectrograms; adding spatial metrics alongside the pooled sliced $W_2$; and investigating the Cartesian geometry's lower phase error at intermediate step counts.
